%======================================================================
% PRÓLOGO
%======================================================================
\chapter*{Prólogo}
\addcontentsline{toc}{chapter}{Prólogo}

El presente documento representa la culminación de años de investigación, desarrollo y aplicación práctica de tecnologías de middleware en el contexto venezolano. Durante mi trayectoria profesional, he tenido la oportunidad de observar de cerca las necesidades tecnológicas de nuestras instituciones públicas y privadas, identificando una brecha significativa entre las soluciones disponibles en el mercado internacional y las realidades específicas de nuestro país.

Las instituciones venezolanas enfrentan desafíos únicos que requieren soluciones igualmente únicas: desde la integración de sistemas heredados desarrollados hace décadas, hasta la necesidad de cumplir con marcos legales específicos como la Ley de Infogobierno y la Ley Especial contra los Delitos Informáticos. Esta realidad nos llevó a concebir Sandra Ecosystem no simplemente como un producto de software, sino como una respuesta integral a las necesidades de transformación digital de nuestras organizaciones.

El ecosistema que se detalla en este documento ha sido diseñado siguiendo principios de Soberanía Tecnológica, priorizando el uso de software libre, estándares abiertos y arquitecturas que permiten el control completo sobre la infraestructura tecnológica. Creemos firmemente que la independencia tecnológica es un pilar fundamental del desarrollo nacional, y este documento busca compartir nuestra visión y avances en esta dirección.

Este trabajo no habría sido posible sin el apoyo de Code Epic Technologies y de todos los profesionales que han contribuido al desarrollo del ecosistema. A ellos mi más sincero agradecimiento.

\vspace{1cm}
\begin{flushright}
\textit{Carlos Peña\\
Caracas, Marzo 2026}
\end{flushright}

\newpage
